\documentclass[11pt]{article}
\usepackage[margin=0.92in]{geometry}
\usepackage{microtype}
\usepackage{amsmath,amssymb}
\usepackage{booktabs}
\usepackage{tabularx}
\usepackage{array}
\usepackage{xcolor}
\usepackage{hyperref}
\usepackage{url}
\usepackage{float}
\hypersetup{colorlinks=true,linkcolor=black,citecolor=black,urlcolor=blue,pdftitle={Supplementary Information for Policy-TD},pdfauthor={Surya Teja Avvaru; Krishna Sai Pokala; Varun Kumar Reddy Dodda},hypertexnames=false}
\newcolumntype{Y}{>{\raggedright\arraybackslash}X}
\newcommand{\policytd}{Policy-TD}
\emergencystretch=3em
\begin{document}

\begin{center}
{\Large\bfseries Supplementary Information for\\Policy-TD: Teacher-Distilled Runtime Control for Frozen Small Language Models\par}
\vspace{0.75em}
\begin{tabular}{@{}>{\centering\arraybackslash}p{0.30\textwidth}>{\centering\arraybackslash}p{0.30\textwidth}>{\centering\arraybackslash}p{0.36\textwidth}@{}}
\textbf{Surya Teja Avvaru} & \textbf{Krishna Sai Pokala} & \textbf{Varun Kumar Reddy Dodda}\\[-0.1em]
{\small\href{mailto:avvarusuryateja@gmail.com}{\texttt{avvarusuryateja@gmail.com}}} &
{\small\href{mailto:ok50250@umbc.edu}{\texttt{ok50250@umbc.edu}}} &
{\small\href{mailto:vdodda1@bowiestate.edu}{\texttt{vdodda1@bowiestate.edu}}}
\end{tabular}
\vspace{0.4em}

{\small Surya Teja Avvaru is an Independent Researcher and the corresponding author.}\\[-0.05em]
{\small Supplementary Information - August 2026}
\end{center}

\section*{S1. Evaluation construction and paired-row accounting}
The internal heldout suite contains 100 prompts across ten families and is evaluated over three guide seeds and three frozen students, producing 900 paired runtime rows per condition. The external-style stress suite contains 300 deterministic prompts and is likewise evaluated over three seeds and three students, producing 2700 rows. The public external suite contains 400 prompts spanning SQuAD2-style answerable QA, SQuAD2-style unanswerable QA, ARC-Challenge, and GSM8K, producing 3600 rows per condition.

The evaluation unit used for helped/harmed accounting is a runtime row: the same prompt, student, and guide setting has both a baseline and a guided outcome. This pairing makes every accuracy change decomposable into recovered baseline failures and corrupted baseline successes. Because each prompt appears across multiple student/seed combinations, runtime rows are not independent prompt draws. The row-level intervals shown in the main paper are therefore descriptive paired intervals, not prompt-cluster generalization intervals.

\section*{S2. Paired metrics and intervention-quality diagnostics}
Let baseline and guided correctness be \(b_i,q_i\in\{0,1\}\). Helped and harmed counts are
\[
H=\sum_i \mathbf{1}[b_i=0\land q_i=1],\qquad
D=\sum_i \mathbf{1}[b_i=1\land q_i=0],
\]
and the paired commitment-control gain in percentage points is
\[
\mathrm{CCG}_{\mathrm{pp}}=100\frac{H-D}{N}.
\]
For behavior-changing interventions, let \(I\) be the number of interventions and \(T\) the number applied to baseline-incorrect rows. We distinguish targeting from outcome:
\[
\mathrm{TargetingPrecision}=\frac{T}{I},\qquad
\mathrm{HelpYield}=\frac{H}{I},\qquad
\mathrm{HarmYield}=\frac{D}{I}.
\]
Targeting precision is a post-hoc diagnostic because baseline correctness is not known at deployment. Help and harm yields quantify whether the chosen bounded action actually changes correctness.

\begin{table}[H]
\centering
\small
\caption{Intervention-quality decomposition from the frozen pooled result tables.}
\begin{tabularx}{\linewidth}{p{0.23\linewidth}p{0.18\linewidth}rrrr}
\toprule
Suite & Setting & \(I\) & Targeting precision & Help yield & Harm yield \\
\midrule
Internal heldout & Conservative & 196 & 78.57\% & 73.98\% & 0.00\% \\
Internal heldout & Higher-recall & 198 & 78.79\% & 74.24\% & 0.00\% \\
External-style stress & Conservative & 190 & 84.21\% & 57.89\% & 0.00\% \\
External-style stress & Higher-recall & 590 & 94.92\% & 18.64\% & 0.00\% \\
Public external & Conservative & 543 & 88.95\% & 4.79\% & 1.10\% \\
Public external & Higher-recall & 572 & 87.94\% & 6.12\% & 2.62\% \\
\bottomrule
\end{tabularx}
\end{table}

The contrast between targeting precision and help yield is a key transfer diagnostic. On the internal heldout, most conservative interventions that target a baseline failure also recover it. On public external evaluation, conservative targeting precision remains 88.95\%, but help yield falls to 4.79\%. Thus public transfer is limited not only by gate selectivity; many correctly targeted failures are not recoverable by the bounded action executor and frozen generator.

\section*{S3. Runtime action semantics}
\begin{table}[H]
\centering
\small
\caption{Runtime action vocabulary and principal failure mode.}
\begin{tabularx}{\linewidth}{p{0.22\linewidth}Y Y}
\toprule
Action & Runtime operation & Principal risk or limitation \\
\midrule
\texttt{finalize} & Accept the current candidate answer. & Preserves a wrong baseline answer if a failure is missed. \\
\texttt{say\_unknown} & Abstain because the prompt lacks sufficient support. & Can harm answerable examples if over-applied. \\
\texttt{repair\_final} & Apply a bounded correction when trace evidence supports a different final answer. & Can over-correct when the trace is unreliable. \\
\texttt{rollback} & Discard an unsafe or unsupported trajectory and request a constrained replacement. & Can discard useful partial reasoning. \\
\texttt{force\_grounded\_step} & Request a constrained continuation grounded in stated evidence. & Cannot create missing generator competence. \\
\texttt{continue} & Permit ordinary continuation without immediate override. & Can delay intervention if the state already warrants action. \\
\bottomrule
\end{tabularx}
\end{table}

\section*{S4. Seed-level runtime stability}
\begin{table}[H]
\centering
\small
\caption{Internal heldout guide-seed results. Each seed is evaluated over three frozen students and 100 heldout prompts.}
\begin{tabular}{lrrrrrr}
\toprule
Seed/setting & Base acc. & Guided acc. & Delta pp & Helped & Harmed & Int. rate \\
\midrule
seed13 conservative & 31.00 & 47.00 & 16.00 & 48 & 0 & 21.67 \\
seed17 conservative & 31.00 & 47.33 & 16.33 & 49 & 0 & 22.00 \\
seed23 conservative & 31.00 & 47.00 & 16.00 & 48 & 0 & 21.67 \\
seed13 higher-recall & 31.00 & 47.33 & 16.33 & 49 & 0 & 22.00 \\
seed17 higher-recall & 31.00 & 47.33 & 16.33 & 49 & 0 & 22.00 \\
seed23 higher-recall & 31.00 & 47.33 & 16.33 & 49 & 0 & 22.00 \\
\bottomrule
\end{tabular}
\end{table}

\section*{S5. Student-level transfer}
Model aliases are Qwen3 for Qwen3-1.7B, SmolLM2 for SmolLM2-1.7B-Instruct, and TinyLlama for TinyLlama-1.1B-Chat.
\begin{table}[H]
\centering
\small
\caption{Student-level behavior on internal heldout and public external evaluation.}
\begin{tabular}{lllrrrrr}
\toprule
Suite & Student & Setting & Base & Guided & Delta pp & Helped & Harmed \\
\midrule
Internal & Qwen3 & Conservative & 49.00 & 61.00 & 12.00 & 36 & 0 \\
Internal & SmolLM2 & Conservative & 19.00 & 35.33 & 16.33 & 49 & 0 \\
Internal & TinyLlama & Conservative & 25.00 & 45.00 & 20.00 & 60 & 0 \\
Internal & Qwen3 & Higher-recall & 49.00 & 61.00 & 12.00 & 36 & 0 \\
Internal & SmolLM2 & Higher-recall & 19.00 & 36.00 & 17.00 & 51 & 0 \\
Internal & TinyLlama & Higher-recall & 25.00 & 45.00 & 20.00 & 60 & 0 \\
Public & Qwen3 & Conservative & 30.75 & 32.67 & 1.92 & 23 & 0 \\
Public & SmolLM2 & Conservative & 8.50 & 8.75 & 0.25 & 3 & 0 \\
Public & TinyLlama & Conservative & 16.75 & 16.25 & -0.50 & 0 & 6 \\
Public & Qwen3 & Higher-recall & 30.75 & 32.67 & 1.92 & 26 & 3 \\
Public & SmolLM2 & Higher-recall & 8.50 & 9.00 & 0.50 & 6 & 0 \\
Public & TinyLlama & Higher-recall & 16.75 & 16.00 & -0.75 & 3 & 12 \\
\bottomrule
\end{tabular}
\end{table}

The student decomposition shows that public transfer depends on the competence and behavior of the frozen generator. Qwen3 transfers positively, SmolLM2 weakly positively, and TinyLlama negatively. The result is consistent with the recoverability interpretation in the main paper.

\section*{S6. Internal family-level mechanism}
\begin{table}[H]
\centering
\small
\caption{Internal heldout family-level effects, pooled over seeds and students. Values are percentages.}
\begin{tabularx}{\linewidth}{Yrrrr}
\toprule
Family & Base & Conservative & Higher-recall & Higher-recall delta pp \\
\midrule
answer-trace consistency & 36.67 & 36.67 & 36.67 & 0.00 \\
conservation & 50.00 & 50.00 & 50.00 & 0.00 \\
entity binding & 66.67 & 66.67 & 66.67 & 0.00 \\
missing information & 13.33 & 97.78 & 100.00 & 86.67 \\
quantifier/negation & 13.33 & 13.33 & 13.33 & 0.00 \\
state transition & 3.33 & 3.33 & 3.33 & 0.00 \\
style pressure & 16.67 & 16.67 & 16.67 & 0.00 \\
target lock & 36.67 & 36.67 & 36.67 & 0.00 \\
unknownability & 23.33 & 100.00 & 100.00 & 76.67 \\
unsupported assumption & 50.00 & 50.00 & 50.00 & 0.00 \\
\bottomrule
\end{tabularx}
\end{table}

\section*{S7. Public external domain decomposition}
\begin{table}[H]
\centering
\small
\caption{Public external results by domain, pooled over seeds and students.}
\begin{tabular}{llrrrrr}
\toprule
Domain & Setting & Base & Guided & Delta pp & Helped & Harmed \\
\midrule
SQuAD2 answerable & Conservative & 28.00 & 29.11 & 1.11 & 10 & 0 \\
SQuAD2 unanswerable & Conservative & 9.33 & 10.00 & 0.67 & 6 & 0 \\
ARC-Challenge & Conservative & 24.33 & 25.11 & 0.78 & 10 & 3 \\
GSM8K math & Conservative & 13.00 & 12.67 & -0.33 & 0 & 3 \\
SQuAD2 answerable & Higher-recall & 28.00 & 28.11 & 0.11 & 10 & 9 \\
SQuAD2 unanswerable & Higher-recall & 9.33 & 11.00 & 1.67 & 15 & 0 \\
ARC-Challenge & Higher-recall & 24.33 & 25.11 & 0.78 & 10 & 3 \\
GSM8K math & Higher-recall & 13.00 & 12.67 & -0.33 & 0 & 3 \\
\bottomrule
\end{tabular}
\end{table}

\section*{S8. Label-transfer diagnostics}
Label-transfer accuracy is diagnostic evidence about guide-head learning, not the primary endpoint. Exact teacher-action imitation can be imperfect even when two actions are behaviorally near-equivalent, and exact action agreement does not guarantee a successful runtime repair if the frozen student cannot complete the requested action.

\begin{table}[H]
\centering
\small
\caption{Released label-transfer diagnostics pooled by student.}
\begin{tabular}{lrrrrrrr}
\toprule
Student & support & answerable & entity & logic & trace & best action & intervene \\
\midrule
SmolLM2 & 48.96 & 42.71 & 78.13 & 73.96 & 32.29 & 36.46 & 50.00 \\
TinyLlama & 26.04 & 56.25 & 64.58 & 62.50 & 66.67 & 20.83 & 56.25 \\
\bottomrule
\end{tabular}
\end{table}

The frozen public label-transfer table contains SmolLM2 and TinyLlama aggregates but no Qwen3 aggregate. This supplement therefore does not infer or fabricate a Qwen3 label-transfer value. Runtime evaluation, which is the primary endpoint, includes all three students.

\section*{S9. Claim-to-artifact map}
\begin{table}[H]
\centering
\small
\caption{Traceability map for the principal manuscript claims.}
\begin{tabularx}{\linewidth}{p{0.34\linewidth}Y}
\toprule
Manuscript claim & Public supporting artifact class \\
\midrule
Internal heldout gain and zero observed harm & frozen pooled runtime tables; seed/student/family decompositions; manifests \\
Router-off returns to baseline & router-off rows in frozen condition tables \\
Mechanistic concentration in missing-information and unknownability families & frozen family-level runtime table \\
External-style positive transfer with zero observed harmed cases & external-style benchmark report, manifest, and frozen tables \\
Public external boundary & public benchmark report and frozen student/domain tables \\
Teacher-label and gate structure & public JSON schema, typed label/action contracts, validation utilities, and diagnostic label-transfer tables \\
Public release integrity & result manifests, checksums, tests, and release-check scripts \\
\bottomrule
\end{tabularx}
\end{table}

\section*{S10. Practical deployment criterion}
For a student-domain pair \((S,D)\), a conservative empirical policy enables guided mode only when paired validation satisfies
\[
\mathrm{CCG}_{\mathrm{pp}}(S,D)\ge 0
\qquad\text{and}\qquad
D_{\mathrm{harm}}(S,D)\le \delta,
\]
where \(\delta\) is the tolerated harmed-case threshold. A strict policy sets \(\delta=0\). This is an empirical admissibility rule, not a correctness or safety guarantee. If the condition is not met, router-off or ordinary baseline behavior is preferred.

\section*{S11. Public release boundary}
The GitHub repository is intentionally curated. It exposes the stable method contract, label schema, runtime action vocabulary, paired metrics, frozen summary evidence, benchmark reports, provenance manifests, tests, and verification utilities. It does not expose raw provider transcripts or credentials, large trace directories, local exploratory scripts, all internal orchestration, or unreleased checkpoints. The public release therefore supports inspection and traceability of aggregate evidence; full from-scratch re-execution of the original research runs requires additional frozen runtime/checkpoint assets beyond the Git repository.

\section*{S12. Author contributions}
Approximate contribution shares for this manuscript are Surya Teja Avvaru 75\%, Krishna Sai Pokala 15\%, and Varun Kumar Reddy Dodda 10\%. Surya Teja Avvaru is the principal contributor, and author order reflects relative contribution.

\end{document}
